%% kapa: exam=1 round=30 year=2019 period=1 form=A subject=civil_law subject_ko=민법 seq=34 number=34
%% source: Q-Net 감정평가사 기출문제 (cst00309) artlSeq=5205205 file=extracted/1008/5205205/2019년 제30회 감정평가사 1차시험 문제지/2019년 제30회 감정평가사 1차시험 문제지/2019년도 제30회 감정평가사 1차 1교시 A형.hwp
%% conversion: kordoc 4.24.51 -> tools/gen1.py (auto); review=proofread-sonnet; flags=none
\begin{question}{34}
지상권에 관한 설명으로 옳지 않은 것은? (다툼이 있으면 판례에 따름)
\begin{choices}
  \item 지상권자는 그 권리의 존속기간 내에서 그 토지를 타인에게 임대할 수 있다.
  \item 구분지상권의 존속기간을 영구적인 것으로 약정하는 것은 허용된다.
  \item 지상권자가 2년 이상의 지료를 지급하지 아니하는 때에는 지상권설정자는 지상권의 소멸을 청구할 수 있다.
  \item 지료연체를 이유로 한 지상권소멸청구에 의해 지상권이 소멸하더라도 지상물매수청구권은 인정된다.
  \item 지상권 설정계약에서 지료의 지급에 대한 약정이 없더라도 지상권의 성립에는 영향이 없다.
\end{choices}
\end{question}
\answer{4}
